\input{../tex-inputs/latex-leseansicht-vorspann}

\section[Gustav Schwarzkopf an Arthur Schnitzler, 8. 7. 1893]{L04281 Gustav Schwarzkopf an Arthur Schnitzler, 8. 7. 1893}
\nopagebreak\mylabel{L04281v}
\rehead{ }\normalsize\beginnumbering\briefempfaengerindex{Schnitzler, Arthur@\textsc{Schnitzler, Arthur}!zzzSchwarzkopf, Gustav@\emph{von Gustav Schwarzkopf}!1893-07-083@{8. 7. 1893}|(be}
\toendnotes[C]{\smallbreak\pagebreak[2]}
\correspDesc{Versand  durch Gustav Schwarzkopf am 8. 7. 1893 in Mödling
\newline{}Erhalt  durch Arthur Schnitzler im Zeitraum [9. 7. 1893 – 13. 7. 1893?] in Bad Ischl}\toendnotes[C]{\smallbreak}
\Standort{CUL, Schnitzler, B 96a.}
\physDesc{Brief, 1 Blatt, 4 Seiten, 1558 Zeichen
\newline{}Handschrift: schwarze Tinte, deutsche Kurrent}\toendnotes[C]{\smallbreak}
\pstart
           \raggedleft{}{\pb}8 VII 93\pend
           
\pstart{}Lieber Freund!\pend\vspace{0.5em}
\pstart
           Beſten Dank für Ihren \label{K_L04281-1v}\edtext{Brief}{\lemma{\textnormal{\emph{Brief}}}\Cendnote{\textnormal{XXXX Auszeichnungsfehler: Dokument L04217 nicht gefunden.}}}\label{K_L04281-1}; es war nett von Ihnen, daß Sie geſchrieben haben.
               Eigentlich habe ich Ihnen nur das zu{ }ſagen, denn mitzuteilen gibt es nichts. Sie
               müſſen das begreiflich finden. We{\geminationn} Sie unter{ }ſo
               günſtigen Bedingungen – \textsc{Ischl\oindex{Bad Ischl@\textbf{Bad Ischl}|pw}}, Ausblick auf die »herrliche Natur« vorzüglich Eigene –
               nichts erleben und noch nicht einmal zur Empfindung reinen Glückes gelangt{ }ſind, wie{ }ſoll ich {\pb}es hier? Sie können{ }ſich
               leicht denken, daß meine So{\geminationm}er Exiſtenz jetzt nicht
               beſonders anregend iſt. Das Einzige, was meinem Leben in den letzten Wochen einen
               Inhalt gab: die neue Köchin\pwindex{?? [Köchin von Gustav Schwarzkopf, 1893] @\textsc{?? [Köchin von Gustav Schwarzkopf, 1893]}|pwv}, bringt
               nun auch keine Überraſchungen mehr. –\pend
           
\pstart
           Von \textsc{Loris\pwindex{Hofmannsthal, Hugo von 1.\,2.\,1874 Wien – 15.\,7.\,1929 Rodaun@\textsc{Hofmannsthal, Hugo von} (1.\,2.\,1874 Wien – 15.\,7.\,1929 Rodaun), \emph{Schriftsteller}|pw}} hab ich geſtern Nachricht beko{\geminationm}en; er iſt mit{ }ſeinen Eltern\pwindex{Hofmannsthal, Hugo August von 21.\,12.\,1841 Wien – 8.\,12.\,1915 ebd.@\textsc{Hofmannsthal, Hugo August von} (21.\,12.\,1841 Wien – 8.\,12.\,1915 ebd.), \emph{Bankdirektor}|pwv}\pwindex{Hofmannsthal, Anna von 27.\,1.\,1849 Wien – 22.\,3.\,1904 Sanatorium Fürth@\textsc{Hofmannsthal, Anna von} (27.\,1.\,1849 Wien – 22.\,3.\,1904 Sanatorium Fürth)|pwv}{ }ſchon wieder in der Fuſch\oindex{Bad Fusch@\textbf{Bad Fusch}|pw},{ }ſie haben es in \textsc{Gossensass\oindex{Gossensaß@\textbf{Gossensaß}, \emph{Hauptstadt}|pw}} nur zwei Tage ausgehalten.Rudolf\pwindex{Schwarzkopf, Rudolf 25.\,5.\,1861 Wien – 13.\,10.\,1893 Meran@\textsc{Schwarzkopf, Rudolf} (25.\,5.\,1861 Wien – 13.\,10.\,1893 Meran), \emph{Schriftsteller}|pw}
               befindet{ }ſich leidlich, er klagt zwar i{\geminationm}er {\pb}und iſt in beſtändiger Angſt, aber wir{ }ſind ja{ }ſchon zufrieden, we{\geminationn}{ }ſeine Angſt{ }ſich als
               grundlos erweiſt. Die nächſte Woche muß er zum großen Teil in Wien\oindex{Wien@\textbf{Wien}, \emph{Verwaltungsgebiet}|pw} zubringen, um{ }ſeine Arbeit zu machen. We{\geminationn} es nur{ }ſchon glücklich vorüber wäre!\pend
           
\pstart
           Meine Brüder\pwindex{Schwarzkopf, Rudolf 25.\,5.\,1861 Wien – 13.\,10.\,1893 Meran@\textsc{Schwarzkopf, Rudolf} (25.\,5.\,1861 Wien – 13.\,10.\,1893 Meran), \emph{Schriftsteller}|pwv}\pwindex{Schwarzkopf, Max 12.\,6.\,1857 Wien – 14.\,4.\,1928 ebd.@\textsc{Schwarzkopf, Max} (12.\,6.\,1857 Wien – 14.\,4.\,1928 ebd.), \emph{Rechtsanwalt}|pwv}\pwindex{Schwarzkopf, Emil 17.\,9.\,1851 Wien – 28.\,1.\,1928 ebd.@\textsc{Schwarzkopf, Emil} (17.\,9.\,1851 Wien – 28.\,1.\,1928 ebd.), \emph{Übersetzer, Komponist}|pwv} danken herzlich für Ihre Grüße und erwidern dieſelben beſtens. Den
               Gruß an Herrn Wachtel\pwindex{Wachtel, Felix Viktor *~15.\,6.\,1859 Lamnitz@\textsc{Wachtel, Felix Viktor} (*~15.\,6.\,1859 Lamnitz), \emph{Regisseur, Schauspieler}|pw} werde ich beſtellen, die
               geforderten »Kühle« werde ich leicht aufbringen. Wo aber{ }ſoll ich die frivole Wärme
               für den Gruß an die Tini\pwindex{Schönberger, Christine 17.\,11.\,1875 – 3.\,2.\,1971 Wien@\textsc{Schönberger, Christine} (17.\,11.\,1875 – 3.\,2.\,1971 Wien), \emph{Gastwirtin}|pw} hernehmen? Da werde
               ich wol vorher einige Scenen aus \textsc{Anatole\pwindex{Schnitzler, Arthur 15.\,5.\,1862 Wien – 21.\,10.\,1931 ebd.@\textsc{Schnitzler, Arthur} (15.\,5.\,1862 Wien – 21.\,10.\,1931 ebd.), \emph{Schriftsteller, Mediziner}!Anatol@\strich\emph{Anatol}|pw}} leſen müſſen {\pb}um mich zu{ }ſtärken
               und in würdiger Weiſe vorzubereiten.\pend
           
\pstart
           \numberlinefalse{}–\numberlinetrue{}\pend
           
\pstart
           Sie verſtändigen mich doch we{\geminationn} Sie wieder in Wien\oindex{Wien@\textbf{Wien}, \emph{Verwaltungsgebiet}|pw}{ }ſind?\pend
           
\pstart
           Herzlichſt{\\[\baselineskip]}Ihr{\\[\baselineskip]}\spacefill\mbox{Gustav Schwarzkopf}\pend
           \leftskip=0em{}
\pstart
           \noindent{}Bitte, beſtellen Sie an D\textsuperscript{r}{ } \textsc{B. Hofma{\geminationn}\pwindex{Beer-Hofmann, Richard 11.\,7.\,1866 Wien – 26.\,9.\,1945 New York City@\textsc{Beer-Hofmann, Richard} (11.\,7.\,1866 Wien – 26.\,9.\,1945 New York City), \emph{Schriftsteller}|pw}} meine allerbeſten Grüße\pend
           \selectlanguage{ngerman}\endnumbering\briefempfaengerindex{Schnitzler, Arthur@\textsc{Schnitzler, Arthur}!zzzSchwarzkopf, Gustav@\emph{von Gustav Schwarzkopf}!1893-07-083@{8. 7. 1893}|)be}\mylabel{L04281h}
\begin{anhang}
\end{anhang}\newcommand{\dateiname}{L04281}\newcommand{\titel}{Gustav Schwarzkopf an Arthur Schnitzler, 8. 7. 1893}\newcommand{\editorInnen}{Herausgegeben von Jahnke, SelmaMüller, Martin Anton}\input{../tex-inputs/latex-leseansicht-abspann}
